\documentclass[10pt,a4paper]{amsart}
\usepackage[margin=19mm]{geometry}
\usepackage{amsmath,amssymb,mathtools}
\usepackage[scr=rsfs,cal=euler]{mathalfa}
\usepackage{xfrac}
\usepackage[unicode,hidelinks,bookmarks=false]{hyperref}
\newcommand{\ie}{i.e.\ }
\newcommand{\cf}{cf.\ }
\newcommand{\resp}{respectively\ }
\newcommand{\Subsection}{\subsection}
\providecommand{\nobreakdash}{\nobreak\hbox{-}}
\setlength{\emergencystretch}{2em}
\multlinegap0pt
\usepackage{float}
\usepackage{amssymb}
\newtheorem{lemma}{Lemma}
\newcommand{\frk}{\mathfrak{k}}
\title{An enhanced Helgason-Johnson bound for $Sp(p, q)$}
\author{Zhan Ying, Chao-Ping Dong}
\date{Source: arXiv:2608.26664v1 (CC BY 4.0)}
\begin{document}
\maketitle

\noindent\emph{Source and licence.} An enhanced Helgason-Johnson bound for $Sp(p, q)$. Version arXiv:2608.26664v1 (CC BY 4.0), distributed by arXiv under the Creative Commons Attribution 4.0 International licence, http://creativecommons.org/licenses/by/4.0/, as recorded in the arXiv licence metadata for this exact version and shown on the abstract page https://arxiv.org/abs/2608.26664v1. Full licence notice: \texttt{licenses/arxiv-2608.26664-CC-BY-4.0.txt}.\par\smallskip

\noindent\textit{Editorial scope.} Counted unit: the article's lemma lem-p (source0.tex lines 1176--1180). The article's complete original proof of it is delivered with it (source0.tex lines 1181--1277, one \texttt{proof} environment of 97 lines). No mathematics is written, repaired or re-derived: the statement and the proof are the article's own lines, in the article's own notation, and no retained line is substituted or re-flowed.  Retained uncounted as the source-local context the proof needs: lines 482--484; lines 638--640.  Nothing was omitted from any retained span, so the package carries no omission record.  No retained line contains a citation, so the wrapper carries no bibliography and no bibliography entry.  No retained line contains a figure environment, an image-inclusion command or drawing code, so no image is delivered and the package declares no asset dependency.  Editorial formatting only, in addition: an AMS \texttt{amsart} preamble that restates the article's own declarations and macros used by the retained lines, namely the environments \texttt{lemma} on separate counters, together with the standard packages those lines need.  The retained environments print in this document as Lemma 1 in order of appearance, whereas the article prints the same items with the numbering of its own full document.  The complete native source of the article as distributed in the arXiv e-print for this exact version is bundled unchanged as \texttt{sources/208655/source0.tex}.
\begin{equation}\label{rho-n-w}
\rho_n^{(w)}=(K_w\mid R_w)=(k_1,\dots,k_p\mid r_1,\dots,r_q).
\end{equation}

 \begin{equation}\label{L-N-M}
L(\mu)=N(\mu)+M(\mu)-\dfrac{e(\mu)}{2}+\sum_{s=1}^n s^2.
\end{equation}

\begin{lemma}\label{lem-p}
Suppose $p\ge 2$. Fix $\rho_n^{(w)}\in\Omega_{p,q}$ as in \eqref{rho-n-w}. Assume in addition that $k_1=q$.
	If for every $0\le s\le q$ and $0\le t\le q$, we have $F_t(\varphi_s)\le B_0$, then
	\[F_t(R_w)=D_w(K_w\mid \phi_t)\le B_0.\]
\end{lemma}

\begin{proof}	For $1\le s\le q$ and $r_{s}\le l\le p-1$, define the $q$-dimensional vector
	\[\varphi_{s,l}=(\underbrace{p-1,\dots,p-1}_{s-1},l,r_{s+1},\dots,r_q)\]
Recall that $R_w=(r_1,\dots,r_q)$.	
Thus $\varphi_{1,r_1}=R_w$. Note also that
\begin{equation}\label{var}
\varphi_{1, p-1}=\varphi_{2, r_2}, \varphi_{2, p-1}=\varphi_{3, r_3}, \dots,
\varphi_{q-1, p-1}=\varphi_{q, r_q}, \varphi_{q, p-1}=\varphi_{q}.
\end{equation}
For $r_{s}+1\le l\le p-1$, set
\[\sigma_{t,s}(l):=F_t(\varphi_{s,l-1})-F_t(\varphi_{s,l}).\]

For $1\le s\le q$ and $0\le l\le p-1$, define
\[\varphi_{s,l}^0=(\underbrace{p-1,\dots,p-1}_{s-1},l,0,\dots,0).\]
For $1\le l\le p-1$, set
	\[\sigma_{t,s}^0(l):=F_t(\varphi^0_{s,l-1})-F_t(\varphi^0_{s,l}).\]
	

We will show that $\sigma_{t,s}(l)$ is monotonically decreasing with respect to $l$ and
	\[\sigma_{t,s}(l)=\sigma_{t,s}^0(l).\]
It follows that for any $r_s\le z\le p-1$,
	\[\sum_{l=z+1}^{p-1}\sigma_{t,s}(l)\le \frac{p-1-z}{p-1}\sum_{l=1}^{p-1}\sigma_{t,s}^0(l).\]
	Note that	\[\sum_{l=z+1}^{p-1}\sigma_{t,s}(l)=F_t(\varphi_{s,z})-F_{t}(\varphi_{s,p-1})\]
	and
\[
\sum_{l=1}^{p-1}\sigma^0_{t,s}(l)
=F_t(\varphi_{s,0}^0)-F_{t}(\varphi_{s,p-1}^0)
=F_t(\varphi_{s-1})-F_t(\varphi_s).\]
	Thus
	\begin{align*}
		F_t(R_w)&= F_t(\varphi_{1, r_1})\\
&=F_t(\varphi_{q}) + \sum_{s=1}^{q} (F_t(\varphi_{s, r_s}) - F_t(\varphi_{s, p-1}))\\
&=F_t(\varphi_{q})+\sum_{s=1}^q\sum_{l=r_s+1}^{p-1}\sigma_{t,s}(l)\\
		&\le F_t(\varphi_{q})+\sum_{s=1}^q  \frac{p-1-r_s}{p-1} \sum_{l=1}^{p-1} \sigma^0_{t,s}(l)\\
&= F_t(\varphi_{q})+\sum_{s=1}^q  \frac{p-1-r_s}{p-1} (F_t(\varphi_{s-1})-F_t(\varphi_s))\\
&=\sum_{s=0}^qc_s F_t(\varphi_s),
	\end{align*}
	where the second step uses \eqref{var} and
	\[c_0=1-\dfrac{r_1}{p-1},\quad c_i=\dfrac{r_{i}-r_{i+1}}{p-1}\quad(1\le i\le q-1),\quad
	c_q=\dfrac{r_q}{p-1}\]
	satisfies $c_s\ge 0$ and $\sum_sc_s=1$. Therefore, the lemma holds.


Now let us analyze $\sigma_{t,s}(l)$, which is the sum of
\begin{equation}\label{term-1}
\|(K(\varphi_{s, l-1})\mid\phi_t) - (K(\varphi_{s, l-1})\mid\varphi_{s, l-1})\|_{\frk}^2 - \|(K(\varphi_{s, l})\mid\phi_t) - (K(\varphi_{s, l})\mid\varphi_{s, l})\|_{\frk}^2
\end{equation}
and
\begin{equation}\label{term-2}
L(K(\varphi_{s, l})\mid\phi_t)- L(K(\varphi_{s, l-1})\mid\phi_t).
\end{equation}

Firstly, we calculate \eqref{term-1}. Note that
\[\|(K(\varphi_{s, l-1})\mid\phi_t) - (K(\varphi_{s, l-1})\mid\varphi_{s, l-1})\|_{\frk}^2=\Phi_q(\phi_t-\varphi_{s,{l-1}})\]
and
\[\|(K(\varphi_{s, l})\mid\phi_t) - (K(\varphi_{s, l})\mid\varphi_{s, l})\|_{\frk}^2=\Phi_q(\phi_t-\varphi_{s,l}).\]
 Suppose $s\le t$. Then
\begin{align*}
	\{\phi_t-\varphi_{s,l}\}_q&=(\underbrace{2p-r_t,\dots,2p-r_{s+1}}_{t-s},\underbrace{2p-l,p+1,\dots,p+1}_{s},r_{t+1},\dots,r_q),\\
	\{\phi_t-\varphi_{s,l-1}\}_q&=(\underbrace{2p-r_t,\dots,2p-r_{s+1}}_{t-s},\underbrace{2p-l+1,p+1,\dots,p+1}_{s},r_{t+1},\dots,r_q).
\end{align*}
Suppose $s>t$. We have
\begin{align*}
	\{\phi_t-\varphi_{s,l}\}_q&=(\underbrace{p+1,\dots,p+1}_{t},\underbrace{p-1,\dots,p-1}_{s-t-1},l,r_{s+1},\dots,r_q),\\
	\{\phi_t-\varphi_{s,l-1}\}_q&=(\underbrace{p+1,\dots,p+1}_{t},\underbrace{p-1,\dots,p-1}_{s-t-1},l-1,r_{s+1},\dots,r_q).
\end{align*}
Thus the value of \eqref{term-1} is
\begin{equation}\label{term-1-1}
\Phi_q(\phi_t-\varphi_{s,{l-1}})-\Phi_q(\phi_t-\varphi_{s,l})=\left\{
\begin{array}{ll}
	-2l+2(n+p+s-t)+1, & s\le t,\\
	-2l-2(q-s)-1,& s>t.
\end{array}
\right.
\end{equation}


	Now compute \eqref{term-2}. Let $(K(\varphi_{s,l})\mid \phi_t)+2\rho_c=(x_1,\dots,x_p\mid y_1,\dots,y_q)$. Note that
	\[K(\varphi_{s,l})-K(\varphi_{s,l-1})=(\underbrace{0,\dots,0,-1}_{p-l+1},0,\dots,0),\]
	and $x_{p-l+1}=q-s+2l$, we have
	\begin{equation}\label{term-2-1}
		N(K(\varphi_{s,l})\mid \phi_t)-N(K(\varphi_{s,l-1})\mid \phi_t)=x_{p-l+1}^2-(x_{p-l+1}+1)^2+2(q+l)=-2l+2s-1.
	\end{equation}
On the other hand,
\begin{equation}\label{term-2-2}
	\begin{split}
	M(K(\varphi_{s,l})\mid \phi_t)-M(K(\varphi_{s,l-1})\mid\phi_t)&=\sum_{j=1}^q(\min\{x_{q-l+1},y_j\}-\min\{x_{q-l+1}+1,y_j\})\\
	&=-2\,\mathrm{card}\{j\mid y_j\ge x_{q-l+1}+1\}\\
	&=-2\, \mathrm{card}\{j\mid y_j\ge q-s+2l+1\}.
	\end{split}
\end{equation}
Finally, $e(K(\varphi_{s,l})\mid \phi_t)-e(K(\varphi_{s,l-1})\mid\phi_t)$ is clearly independent of $r_{s+1},\dots,r_q$. Since \eqref{term-1} and \eqref{term-2} remain unchanged when $r_{s+1}=\cdots=r_q=0$, then
\[F_t(\varphi_{s,l-1})-F_t(\varphi_{s,l})=F_t(\varphi_{s,l-1}^0)-F_t(\varphi_{s,l}^0).\]

Note that $|e(K(\varphi_{s,l})\mid \phi_t)-e(K(\varphi_{s,l-1})\mid \phi_t)|\le 1$, then by \eqref{L-N-M}, \eqref{term-1-1}, \eqref{term-2-1}, \eqref{term-2-2},
\[\sigma_{t,s}(l)-\sigma_{t,s}(l+1)\ge 4-2-1=1>0.\]
This shows that $\sigma_{t,s}(l)$ decreases as $l$ increases and we done.
\end{proof}

\end{document}
